\subsection{Terminal effects across the current factorials and weighting ablation}
\label{app:current-campaign-results}

Tables~\ref{tab:current-e118}--\ref{tab:current-e120} collect every registered
model--domain contrast in the September 9 endpoint snapshot, including
comparisons whose full five-seed block is still accumulating. Each row uses
the exact terminal seed intersection for that contrast. We report correctness
$P=\texttt{pass@8}$, raw support $D=\texttt{distinct@8}$, and extra correct
modes $B=D-P$. The last quantity counts sampled alternatives beyond the first
correct outcome; it remains coupled to correctness. A dash denotes no observed
terminal pair. Partial means are descriptive, with their observed $n/5$ and
no interval. Completed-block intervals are nominal 95\% estimates, without
multiple-comparison or repeated-look adjustment.

\begin{table}[H]
  \centering
  \caption{\textbf{ReplayMaxRL effects across model scales at pass 8.}
  All effects are ReplayMaxRL minus MaxRL on the displayed paired seeds.
  Brackets give paired Student-$t$ intervals for $\Delta B$ in complete blocks.
  The machine-readable report also retains both arm means, all per-seed
  effects, and intervals for $P$ and $D$.}
  \label{tab:current-e118}
  \scriptsize
  \setlength{\tabcolsep}{3pt}
  \resizebox{\linewidth}{!}{%
  \begin{tabular}{@{}lllcrrr@{}}
    \toprule
    \textbf{Model} & \textbf{Domain} & \textbf{Objective} & $n/5$ &
    $\Delta P$ & $\Delta D$ & $\Delta B$ [95\% interval] \\
    \midrule
    % Model & Domain & Contrast & n/5 & Delta P8 & Delta D8 & Delta B8 [95%]
% E118/E119: replay minus base. E120: uniform minus frequency.
% Intervals are stored, descriptive, and only present at n=5.
% E120 dagger: at least one available endpoint pair lacks mechanism eligibility.
Qwen-0.5B & Graph & Re:MaxRL & 5/5 & $+0.408$ & $+1.630$ & $+1.221\;[+0.926, +1.517]$ \\
Qwen-0.5B & Countdown & Re:MaxRL & 5/5 & $+0.084$ & $+1.074$ & $+0.990\;[+0.859, +1.121]$ \\
Qwen-0.5B & Python & Re:MaxRL & 5/5 & $+0.509$ & $+0.955$ & $+0.446\;[-0.156, +1.048]$ \\
Qwen-0.5B & MathIR & Re:MaxRL & 5/5 & $+0.343$ & $+0.356$ & $+0.013\;[+0.008, +0.017]$ \\
Qwen-0.5B & PantryPlan & Re:MaxRL & 5/5 & $+0.189$ & $+0.941$ & $+0.751\;[+0.520, +0.982]$ \\
Falcon-1B & Graph & Re:MaxRL & 5/5 & $+0.014$ & $-0.093$ & $-0.107\;[-0.220, +0.005]$ \\
Falcon-1B & Countdown & Re:MaxRL & 5/5 & $+0.025$ & $+0.133$ & $+0.108\;[+0.082, +0.134]$ \\
Falcon-1B & Python & Re:MaxRL & 5/5 & $+0.127$ & $+0.280$ & $+0.153\;[-0.272, +0.578]$ \\
Falcon-1B & MathIR & Re:MaxRL & 5/5 & $+0.424$ & $+0.443$ & $+0.019\;[+0.006, +0.031]$ \\
Falcon-1B & PantryPlan & Re:MaxRL & 5/5 & $+0.073$ & $+0.375$ & $+0.302\;[+0.034, +0.569]$ \\
Qwen-3B & Graph & Re:MaxRL & 4/5 & $+0.042$ & $+0.231$ & $+0.189$ \\
Qwen-3B & Countdown & Re:MaxRL & 0/5 & --- & --- & --- \\
Qwen-3B & Python & Re:MaxRL & 5/5 & $+0.277$ & $+0.391$ & $+0.115\;[+0.080, +0.149]$ \\
Qwen-3B & MathIR & Re:MaxRL & 1/5 & $+0.234$ & $+0.242$ & $+0.008$ \\
Qwen-3B & PantryPlan & Re:MaxRL & 1/5 & $+0.039$ & $+0.613$ & $+0.574$ \\
\bottomrule

  \end{tabular}}
\end{table}

At Qwen2.5-0.5B, ReplayMaxRL adds $1.221$, $.990$, and $.751$ extra
correct modes on Graph, Countdown, and PantryPlan, respectively. Each
five-seed interval excludes zero. At Falcon3-1B, the corresponding breadth
effect is smaller and Graph's mean is negative with an interval spanning
zero. Thus the positive aggregate correctness and raw-support effects do
not imply a breadth gain in every domain. The completed Qwen2.5-3B Python
block and the partial scale extensions remain separate comparisons.

\begin{table}[H]
  \centering
  \caption{\textbf{Replay effects on the harder Level-2 tasks at pass 8.}
  MaxRL and Dr.GRPO identify the fresh objective; each effect is its replay
  variant minus the corresponding objective without replay. Each row uses
  its own pairwise intersection, which can exceed the four-method
  intersection. Brackets give paired Student-$t$ intervals for $\Delta B$
  only when all five registered pairs are available.}
  \label{tab:current-e119}
  \scriptsize
  \setlength{\tabcolsep}{3pt}
  \resizebox{\linewidth}{!}{%
  \begin{tabular}{@{}lllcrrr@{}}
    \toprule
    \textbf{Model} & \textbf{Domain} & \textbf{Objective} & $n/5$ &
    $\Delta P$ & $\Delta D$ & $\Delta B$ [95\% interval] \\
    \midrule
    % Model & Domain & Contrast & n/5 & Delta P8 & Delta D8 & Delta B8 [95%]
% E118/E119: replay minus base. E120: uniform minus frequency.
% Intervals are stored, descriptive, and only present at n=5.
% E120 dagger: at least one available endpoint pair lacks mechanism eligibility.
Qwen-0.5B & Graph & Dr.GRPO & 5/5 & $+0.561$ & $+1.019$ & $+0.458\;[+0.394, +0.522]$ \\
Qwen-0.5B & Graph & MaxRL & 5/5 & $+0.305$ & $+0.668$ & $+0.363\;[+0.280, +0.446]$ \\
Qwen-0.5B & Countdown & Dr.GRPO & 5/5 & $+0.025$ & $+0.045$ & $+0.020\;[-0.014, +0.055]$ \\
Qwen-0.5B & Countdown & MaxRL & 5/5 & $+0.040$ & $+0.057$ & $+0.017\;[+0.001, +0.032]$ \\
Qwen-0.5B & Python & Dr.GRPO & 5/5 & $+0.075$ & $+0.400$ & $+0.325\;[-0.228, +0.878]$ \\
Qwen-0.5B & Python & MaxRL & 5/5 & $+0.037$ & $+0.212$ & $+0.175\;[-0.295, +0.645]$ \\
Qwen-0.5B & MathIR & Dr.GRPO & 5/5 & $+0.493$ & $+0.493$ & $+0.000\;[+0.000, +0.000]$ \\
Qwen-0.5B & MathIR & MaxRL & 5/5 & $+0.173$ & $+0.175$ & $+0.002\;[-0.003, +0.006]$ \\
Qwen-0.5B & PantryPlan & Dr.GRPO & 0/5 & --- & --- & --- \\
Qwen-0.5B & PantryPlan & MaxRL & 0/5 & --- & --- & --- \\
\bottomrule

  \end{tabular}}
\end{table}

Graph provides the clearest completed Level-2 evidence: replay increases
both correctness and extra modes under both fresh objectives. Python and
MathIR have positive mean correctness and raw-support effects, but their
intervals include zero. MathIR's extra-mode effects are close to zero,
illustrating the distinction between recovering correctness and producing
multiple correct outcomes.

The four-arm analysis also separates the estimator effect from its
interaction with replay (Table~\ref{tab:level2-factorial-contrasts}). On
Graph, MaxRL without replay raises $P$ by $+.233$ relative to Dr.GRPO
($[+.140,+.325]$). Replay remains beneficial under both estimators, but its
marginal gain is smaller under MaxRL: the interaction is $-.256$ on $P$
($[-.329,-.183]$), $-.351$ on $D$ ($[-.446,-.255]$), and $-.095$ on $B$
($[-.186,-.003]$). This supports an additional replay benefit after stronger
fresh optimization, without a superadditive interaction.

\begin{table}[H]
  \centering
  \caption{\textbf{Fresh-objective effects and replay interactions on Level 2.}
  All three contrasts within a domain use the same four-arm seed
  intersection. The interaction is
  $(\mathrm{ReplayMaxRL}-\mathrm{MaxRL})-
  (\mathrm{ReplayDr.GRPO}-\mathrm{Dr.GRPO})$.
  Each entry gives its mean and, only for a complete five-seed block, a
  paired unadjusted 95\% Student-$t$ interval.}
  \label{tab:level2-factorial-contrasts}
  \scriptsize
  \setlength{\tabcolsep}{3pt}
  \resizebox{\linewidth}{!}{%
  \begin{tabular}{@{}lclrrr@{}}
    \toprule
    \textbf{Domain} & $n$ & \textbf{Contrast} &
    $\Delta P$ & $\Delta D$ & $\Delta B$ \\
    \midrule
    % Domain & paired n & contrast & delta pass@8 & delta distinct@8 & delta extra modes
Graph & 5 & MaxRL $-$ Dr.GRPO & $+.233\;[+.140,+.325]$ & $+.322\;[+.162,+.482]$ & $+.089\;[+.019,+.159]$ \\
Graph & 5 & Re:MaxRL $-$ Re:Dr.GRPO & $-.023\;[-.062,+.016]$ & $-.029\;[-.198,+.140]$ & $-.006\;[-.142,+.130]$ \\
Graph & 5 & Replay $\times$ MaxRL interaction & $-.256\;[-.329,-.183]$ & $-.351\;[-.446,-.255]$ & $-.095\;[-.186,-.003]$ \\
Countdown & 5 & MaxRL $-$ Dr.GRPO & $-.014\;[-.073,+.045]$ & $-.014\;[-.073,+.045]$ & $+.000\;[+.000,+.000]$ \\
Countdown & 5 & Re:MaxRL $-$ Re:Dr.GRPO & $+.002\;[-.008,+.012]$ & $-.002\;[-.047,+.044]$ & $-.004\;[-.042,+.035]$ \\
Countdown & 5 & Replay $\times$ MaxRL interaction & $+.016\;[-.047,+.078]$ & $+.012\;[-.068,+.092]$ & $-.004\;[-.042,+.035]$ \\
Python & 5 & MaxRL $-$ Dr.GRPO & $+.000\;[+.000,+.000]$ & $+.000\;[+.000,+.000]$ & $+.000\;[+.000,+.000]$ \\
Python & 5 & Re:MaxRL $-$ Re:Dr.GRPO & $-.037\;[-.232,+.157]$ & $-.188\;[-.988,+.613]$ & $-.150\;[-1.014,+.713]$ \\
Python & 5 & Replay $\times$ MaxRL interaction & $-.037\;[-.232,+.157]$ & $-.188\;[-.988,+.613]$ & $-.150\;[-1.014,+.713]$ \\
MathIR & 5 & MaxRL $-$ Dr.GRPO & $+.283\;[-.193,+.759]$ & $+.283\;[-.193,+.759]$ & $+.000\;[+.000,+.000]$ \\
MathIR & 5 & Re:MaxRL $-$ Re:Dr.GRPO & $-.036\;[-.158,+.085]$ & $-.035\;[-.158,+.088]$ & $+.002\;[-.003,+.006]$ \\
MathIR & 5 & Replay $\times$ MaxRL interaction & $-.320\;[-.864,+.225]$ & $-.318\;[-.860,+.224]$ & $+.002\;[-.003,+.006]$ \\
PantryPlan & 0 & MaxRL $-$ Dr.GRPO & \textemdash & \textemdash & \textemdash \\
PantryPlan & 0 & Re:MaxRL $-$ Re:Dr.GRPO & \textemdash & \textemdash & \textemdash \\
PantryPlan & 0 & Replay $\times$ MaxRL interaction & \textemdash & \textemdash & \textemdash \\
    \bottomrule

  \end{tabular}}
\end{table}

\begin{table}[H]
  \centering
  \caption{\textbf{Uniform versus fresh-frequency replay at pass 8.}
  Effects are uniform minus frequency weighting. Brackets give the paired
  percentile bootstrap interval for $\Delta B$ over all $5^5$ ordered
  resamples in complete blocks. Every observed treatment endpoint in this
  snapshot passes the persisted replay-weight telemetry audit. Partial
  blocks remain descriptive; Qwen2.5-3B PantryPlan has no terminal pair.}
  \label{tab:current-e120}
  \scriptsize
  \setlength{\tabcolsep}{3pt}
  \resizebox{\linewidth}{!}{%
  \begin{tabular}{@{}lllcrrr@{}}
    \toprule
    \textbf{Model} & \textbf{Domain} & \textbf{Weighting} & $n/5$ &
    $\Delta P$ & $\Delta D$ & $\Delta B$ [95\% interval] \\
    \midrule
    % Model & Domain & Contrast & n/5 & Delta P8 & Delta D8 & Delta B8 [95%]
% E118/E119: replay minus base. E120: uniform minus frequency.
% Intervals are stored, descriptive, and only present at n=5.
% E120 dagger: at least one available endpoint pair lacks mechanism eligibility.
Qwen-0.5B & Graph & Uniform & 5/5 & $+0.066$ & $+0.662$ & $+0.596\;[+0.513, +0.687]$ \\
Qwen-0.5B & Countdown & Uniform & 5/5 & $+0.022$ & $+0.661$ & $+0.639\;[+0.544, +0.733]$ \\
Qwen-0.5B & Python & Uniform & 5/5 & $-0.077$ & $-0.040$ & $+0.037\;[-0.012, +0.123]$ \\
Qwen-0.5B & MathIR & Uniform & 5/5 & $-0.005$ & $+0.014$ & $+0.019\;[-0.001, +0.038]$ \\
Qwen-0.5B & PantryPlan & Uniform & 5/5 & $+0.043$ & $+0.341$ & $+0.298\;[+0.174, +0.400]$ \\
Falcon-1B & Graph & Uniform & 5/5 & $+0.045$ & $+0.278$ & $+0.233\;[+0.182, +0.320]$ \\
Falcon-1B & PantryPlan & Uniform & 3/5 & $-0.058$ & $-0.598$ & $-0.540$ \\
Qwen-3B & Graph & Uniform & 2/5 & $+0.004$ & $+0.142$ & $+0.138$ \\
Qwen-3B & PantryPlan & Uniform & 0/5 & --- & --- & --- \\
\bottomrule

  \end{tabular}}
\end{table}

The complete Graph, Countdown, and PantryPlan blocks at Qwen2.5-0.5B
favor uniform weighting on breadth, as does Falcon3-1B Graph. The
Qwen2.5-3B Graph extension also favors uniform weighting in its two observed
pairs ($\Delta B=+.138$), whereas Falcon3-1B PantryPlan favors frequency
weighting in its three observed pairs ($\Delta B=-.540$). Neither partial
block receives an interval. These differences argue against a universal
weighting advantage. The registered Qwen2.5-0.5B five-domain mean remains
$\Delta B=+.318$ with interval $[+.272,+.358]$; its correctness interval
$[-.119,+.139]$ establishes neither a gain nor equivalence.

The deterministic generator
\path{ops/exp_scaling/build_paper_current_campaign_results.py} consumes the
dated census and writes these tables, a readable report, and CSV/JSON
containing full-precision paired arm means, effects, seed sets, and
uncertainty. The original September 4 weighting analysis is preserved.
